\documentclass[11pt,a4paper]{article}
\usepackage[utf8]{inputenc}
\usepackage[T1]{fontenc}
\usepackage{amsmath,amssymb,amsthm}
\usepackage{geometry}
\usepackage{booktabs}
\usepackage{array}
\usepackage{xcolor}
\usepackage{hyperref}
\geometry{margin=2.3cm}
\hypersetup{colorlinks=true, linkcolor=blue!60!black, urlcolor=blue!60!black}
\newcolumntype{L}[1]{>{\raggedright\arraybackslash}p{#1}}

\newtheorem*{claim}{Claim}

\title{\textbf{Proof} \\ \large Steps 2--3: from the discrete objective\\
to the continuum integral, HJB equation,\\and matrix Riccati ODE}
\author{Full derivation --- see WorkPlan.tex Stage~2\\
and Summary.tex for the one-line status}
\date{}

\begin{document}
\maketitle

\section*{Notation}

\begin{center}
\begin{tabular}{@{}L{2.1cm}L{12.5cm}@{}}
\toprule
\textbf{Symbol} & \textbf{Meaning} \\
\midrule
$M$ & Number of assets. \\
$N$, $\tau$ & Number of discrete time steps, and the step size ($\tau=T/N$), in Step~1's discretisation. \\
$j$, $k$ & Asset index ($j=1,\dots,M$) and time-step index ($k=1,\dots,N$). \\
$\gamma_j$ & Permanent-impact coefficient for asset $j$. \\
$\varepsilon_j$ & Spread-cost coefficient for asset $j$. \\
$\tilde\eta_j$ & Temporary-impact coefficient for asset $j$; $\widetilde H_\tau=\mathrm{diag}(\tilde\eta_j)$. \\
$n_{jk}$ & Shares of asset $j$ traded at step $k$ (discrete). \\
$\mathbf x_k$, $\mathbf x(t)$ & Inventory vector at step $k$ (discrete) / time $t$ (continuous). \\
$\mathbf v(t)$ & Selling-rate vector, $\mathbf v(t)=-\dot{\mathbf x}(t)$. \\
$\mathbf X$ & Initial inventory vector, $\mathbf x(0)=\mathbf X$. \\
$T$ & Liquidation horizon. \\
$\lambda$ & Risk-aversion parameter. \\
$C$ & Risk/covariance matrix ($M\times M$, symmetric, positive semi-definite). \\
$H$ & Continuum temporary-impact matrix $\mathrm{diag}(\eta_j)\succ0$; $\widetilde H_\tau\to H$ as $\tau\to0$. \\
$J[\mathbf v]$ & The continuum objective, \eqref{eq:J}, to be minimised over schedules $\mathbf v(\cdot)$. \\
$\Phi$ & Terminal penalty matrix; $\mathcal V(\mathbf x,T)=\tfrac12\mathbf x^\top\Phi\mathbf x$ enforces $\mathbf x(T)=\mathbf 0$ as $\Phi\to\infty$. \\
$\mathcal V(\mathbf x,t)$ & Cost-to-go (value function) from state $\mathbf x$ at time $t$. \\
$P(t)$ & Symmetric matrix with $\mathcal V(\mathbf x,t)=\tfrac12\mathbf x^\top P(t)\mathbf x$; solves the Riccati ODE \eqref{eq:ric}. \\
$\mathbf v^*$ & Optimal feedback control, $\mathbf v^*=\tfrac12H^{-1}P(t)\mathbf x=\tfrac12H^{-1}\nabla_{\mathbf x}\mathcal V$. \\
$U$, $\mu_i$ & Orthogonal eigenbasis and eigenvalues of $H^{-1/2}CH^{-1/2}=U\,\mathrm{diag}(\mu_i)\,U^\top$. \\
$\kappa_i$ & Decay rate of decoupled mode $i$, $\kappa_i=\sqrt{\lambda\mu_i}$. \\
$Q(t)$, $\hat Q(t)$, $q_i(t)$ & Rescaled ($H^{-1/2}PH^{-1/2}$), rotated ($U^\top QU$), and diagonal-entry forms of $P(t)$ used in the Corollary's proof. \\
$\phi(t)$ & Riccati-linearising substitution, $q=-2\dot\phi/\phi$, used to solve the scalar Riccati ODE in closed form. \\
$\mathbf z(t)$, $z_i(t)$, $K_i$ & Decoupled coordinates $\mathbf z=U^\top H^{1/2}\mathbf x$, their components, and the integration constants fixing $z_i(0)$. \\
\bottomrule
\end{tabular}
\end{center}

\section*{Step 2 --- from the discrete objective to the continuum integral}

Starting point, Step~1's discrete objective (Summary.tex Step~1):
\begin{equation}
U[\mathbf x]=\sum_{j=1}^{M}\Big[\tfrac12\gamma_jX_j^2+\varepsilon_j\sum_{k=1}^N
|n_{jk}|+\frac{\tilde\eta_j}{\tau}\sum_{k=1}^N n_{jk}^2\Big]+\lambda\sum_{k=1}^N
\tau\,\mathbf x_k^\top C\mathbf x_k .
\label{eq:s1}
\end{equation}

\textbf{1. Separate the fixed offset.} Under sales-only trading and exact
liquidation, the two schedule-independent terms have fixed value
\[
K(\mathbf X)=\sum_j\left(\tfrac12\gamma_jX_j^2+\varepsilon_jX_j\right).
\]
The spread term is not constant if a finite terminal penalty permits residual
inventory. For exact liquidation, $K(\mathbf X)$ is kept in the total cost and
omitted only from the optimization below. $\tfrac12\gamma_jX_j^2$
(permanent impact) depends only on the total amount liquidated, not on how it
is spread over time, so it is a constant over all schedules. Under the
sales-only assumption A7 ($n_{jk}\ge0$ always), $|n_{jk}|=n_{jk}$ and
$\sum_kn_{jk}=X_j$, so $\varepsilon_j\sum_k|n_{jk}|=\varepsilon_jX_j$ is also a
constant. Neither term can be changed by choosing a different schedule, so
the schedule-dependent objective is
\[
\tilde U[\mathbf x]=\sum_{j}\frac{\tilde\eta_j}{\tau}\sum_kn_{jk}^2
+\lambda\sum_k\tau\,\mathbf x_k^\top C\mathbf x_k .
\]

\textbf{2. Substitute $n_{jk}=\tau v_j(t_k)$} (trade size = rate $\times$
time-step):
\[
\frac{\tilde\eta_j}{\tau}\sum_kn_{jk}^2=\frac{\tilde\eta_j}{\tau}\sum_k\tau^2
v_j(t_k)^2=\sum_k\tau\,\tilde\eta_j\,v_j(t_k)^2 .
\]
At a fixed time step, stacking over $j$ with $\widetilde H_\tau=\mathrm{diag}(\tilde\eta_j)$:
\[
\tilde U=\sum_k\tau\Big(\mathbf v(t_k)^\top H\,\mathbf v(t_k)+\lambda\,\mathbf
x(t_k)^\top C\,\mathbf x(t_k)\Big),
\]
a Riemann sum with step size $\tau$ (using $H=\widetilde H_\tau$
in this discrete display). In the continuum equations below,
$H=\lim_{\tau\to0}\widetilde H_\tau=\operatorname{diag}(\eta_j)$.

\textbf{3. Take the continuum limit $\tau\to0$} ($N=T/\tau\to\infty$). A
Riemann sum $\sum_k\tau f(t_k)\to\int_0^Tf(t)\,dt$ for bounded, (piecewise)
continuous $f$, so
\begin{equation}
J[\mathbf v]=\int_0^T\big(\mathbf v(t)^\top H\mathbf v(t)+\lambda\,\mathbf
x(t)^\top C\mathbf x(t)\big)\,dt .
\label{eq:J}
\end{equation}

\textbf{4. The inventory recursion becomes the state equation.}
$\mathbf x_k=\mathbf x_{k-1}-\mathbf n_k=\mathbf x_{k-1}-\tau\,\mathbf v(t_k)$,
so $(\mathbf x_k-\mathbf x_{k-1})/\tau=-\mathbf v(t_k)$; as $\tau\to0$ this is
the definition of the derivative,
\begin{equation}
\dot{\mathbf x}=-\mathbf v,\qquad \mathbf x(0)=\mathbf X,\qquad
\mathbf x(T)=\mathbf 0,
\label{eq:state}
\end{equation}
the boundary conditions carrying over directly from the discrete recursion's
endpoints.

Thus, for exact liquidation and sales-only trading,
\begin{equation}
J_{\mathrm{full}}[\mathbf v]=K(\mathbf X)+J[\mathbf v].
\end{equation}
The HJB and Riccati equations below optimize $J$; the full optimal initial cost
is $K(\mathbf X)+\mathcal V(\mathbf X,0)$. A finite-$\Phi$ approximation that
allows nonzero terminal inventory does not preserve the spread offset as a
constant.

This is the multi-asset continuum limit of
Bergault--Drissi--Gu\'eant~\cite{bergault2022}, Sec.~2 and App.~A; the
single-asset case matches Almgren--Chriss~\cite{almgrenchriss2000},
Sec.~2.2 and App.~A.2.

\section*{Step 3 --- what is being proved next}

Starting from Step~2's schedule-dependent continuum objective
\eqref{eq:J}--\eqref{eq:state},
with $H$ diagonal positive-definite and $C$ symmetric positive semi-definite
(assumptions A2/A11, A13), this note derives, by Bellman's dynamic
programming principle (not Euler--Lagrange, since a state-space PDE is what
Stage~7's Schr\"odingerisation needs downstream), exactly the two boxed
results of Summary.tex Step~3:
\begin{equation}
\partial_t\mathcal V-\tfrac14(\nabla_{\mathbf x}\mathcal V)^\top H^{-1}
(\nabla_{\mathbf x}\mathcal V)+\lambda\,\mathbf x^\top C\mathbf x=0,
\label{eq:hjb}
\end{equation}
\begin{equation}
\dot P=\tfrac12PH^{-1}P-2\lambda C,\qquad P(T)=\Phi.
\label{eq:ric}
\end{equation}
A corollary at the end also recovers the closed-form trajectory used in
Stage~2--3, for the hard-liquidation limit $\Phi\to\infty$.

\section*{Step A --- Bellman's principle gives the HJB PDE}

Define the cost-to-go from state $\mathbf x$ at time $t$:
\[
\mathcal V_\Phi(\mathbf x,t)=\inf_{\mathbf v(\cdot)}\left\{\int_t^T\big(\mathbf
v(s)^\top H\mathbf v(s)+\lambda\,\mathbf x(s)^\top C\mathbf x(s)\big)\,ds
+\tfrac12\mathbf x(T)^\top\Phi\mathbf x(T)\right\},
\qquad \dot{\mathbf x}=-\mathbf v,\ \ \mathbf x(t)=\mathbf x .
\]
Here the controls are unconstrained real vectors and $\Phi\succeq0$.
Exact liquidation instead restricts paths to $\mathbf x(T)=0$ and uses
singular terminal data. For sales-only controls, the minimum must be over
$\mathbf v\ge0$; the quadratic feedback below solves that problem at an
initial state only if its resulting path is feasible. The all-time check
for the chosen portfolio is in WorkPlan.tex Stage~3.

Split the horizon into $[t,t+dt]$ and $[t+dt,T]$ and use that the second
piece, started from whatever state the first piece reaches, is again optimal
(Bellman's principle):
\[
\mathcal V(\mathbf x,t)=\min_{\mathbf v}\Big[\big(\mathbf v^\top H\mathbf
v+\lambda\,\mathbf x^\top C\mathbf x\big)\,dt+\mathcal V(\mathbf x-\mathbf
v\,dt,\,t+dt)\Big]+o(dt).
\]
Taylor-expand the second term to first order in $dt$:
\[
\mathcal V(\mathbf x-\mathbf v\,dt,\,t+dt)=\mathcal V(\mathbf x,t)+\partial_t
\mathcal V\,dt-\mathbf v^\top\nabla_{\mathbf x}\mathcal V\,dt+o(dt).
\]
Substitute back, cancel $\mathcal V(\mathbf x,t)$ from both sides, divide by
$dt$, and let $dt\to0$:
\begin{equation}
0=\min_{\mathbf v}\Big[\mathbf v^\top H\mathbf v-\mathbf
v^\top\nabla_{\mathbf x}\mathcal V\Big]+\partial_t\mathcal V+\lambda\,\mathbf
x^\top C\mathbf x.
\label{eq:pre-hjb}
\end{equation}

\section*{Step B --- pointwise minimization: the optimal feedback control}

The bracket in \eqref{eq:pre-hjb} is a convex quadratic in $\mathbf v$ (since
$H\succ0$). Setting its gradient to zero,
\[
2H\mathbf v-\nabla_{\mathbf x}\mathcal V=0\quad\Longrightarrow\quad
\mathbf v^*=\tfrac12H^{-1}\nabla_{\mathbf x}\mathcal V.
\]
Substituting $\mathbf v^*$ back into the bracket:
\[
\mathbf v^{*\top}H\mathbf v^*-\mathbf v^{*\top}\nabla_{\mathbf x}\mathcal V
=\tfrac14(\nabla_{\mathbf x}\mathcal V)^\top H^{-1}(\nabla_{\mathbf x}\mathcal
V)-\tfrac12(\nabla_{\mathbf x}\mathcal V)^\top H^{-1}(\nabla_{\mathbf x}
\mathcal V)=-\tfrac14(\nabla_{\mathbf x}\mathcal V)^\top
H^{-1}(\nabla_{\mathbf x}\mathcal V).
\]
Plugging this into \eqref{eq:pre-hjb} gives exactly \eqref{eq:hjb}, with
terminal condition $\mathcal V(\mathbf x,T)=\tfrac12\mathbf x^\top\Phi\mathbf
x$ (the penalty form of the hard constraint $\mathbf x(T)=\mathbf 0$, A10).
This matches Almgren--Chriss~\cite{almgrenchriss2000}, App.~A.2, and
Bergault--Drissi--Gu\'eant~\cite{bergault2022}, Prop.~1 (their Eqs.~16--17,
under the identification $\eta=H$, $\Sigma=C$, $\gamma=2\lambda$, no
Ornstein--Uhlenbeck signal, i.e.\ $R=0$).

\section*{Step C --- the quadratic ansatz collapses the PDE to a Riccati ODE}

Try $\mathcal V(\mathbf x,t)=\tfrac12\mathbf x^\top P(t)\mathbf x+c(t)$ with
$P(t)$ symmetric. Then $\partial_t\mathcal V=\tfrac12\mathbf x^\top\dot
P\mathbf x+\dot c$ and $\nabla_{\mathbf x}\mathcal V=P(t)\mathbf x$. Substituting
into \eqref{eq:hjb}:
\[
\tfrac12\mathbf x^\top\dot P\mathbf x-\tfrac14\mathbf x^\top PH^{-1}P\mathbf
x+\lambda\,\mathbf x^\top C\mathbf x+\dot c=0\qquad\text{for all }\mathbf x.
\]
At $\mathbf x=0$ this gives $\dot c=0$; $c(T)=0$ implies $c\equiv0$.
A quadratic form $\mathbf x^\top M\mathbf x$ vanishes for all $\mathbf x$ iff
the symmetric part of $M$ is zero; since $\dot P$, $PH^{-1}P$ and $C$ are
each already symmetric, this forces
\[
\tfrac12\dot P-\tfrac14PH^{-1}P+\lambda C=0\quad\Longrightarrow\quad
\dot P=\tfrac12PH^{-1}P-2\lambda C,
\]
which is exactly \eqref{eq:ric}, with $P(T)=\Phi$ read off the terminal
condition above. The optimal control becomes the linear feedback law
$\mathbf v^*(\mathbf x,t)=\tfrac12H^{-1}P(t)\mathbf x$. This matches
Bergault--Drissi--Gu\'eant~\cite{bergault2022}, Prop.~2 and App.~A
(their Eqs.~20--21, same identification as above, plus $\Gamma=\tfrac12\Phi$,
$P_{\text{here}}=-2A_{\text{theirs}}$); the one-asset scalar case matches
Almgren--Chriss~\cite{almgrenchriss2000}, Eq.~17.

\section*{Corollary --- closed-form trajectory, $\Phi\to\infty$}

\begin{claim}
Let $H^{-1/2}CH^{-1/2}=U\,\mathrm{diag}(\mu_i)\,U^\top$ ($U$ orthogonal,
$\mu_i\ge0$ since $C\succeq0$) and $\kappa_i=\sqrt{\lambda\mu_i}$. As
$\Phi\to\infty$ (hard liquidation, $\mathbf x(T)=\mathbf 0$),
\[
\mathbf x(t)=H^{-1/2}U\,\mathrm{diag}\!\Big(\frac{\sinh\kappa_i(T-t)}
{\sinh\kappa_iT}\Big)U^\top H^{1/2}\mathbf X .
\]
\end{claim}

\begin{proof}
Write $Q(t)=H^{-1/2}P(t)H^{-1/2}$ (symmetric). Using $H^{-1}=H^{-1/2}H^{-1/2}$,
\[
\dot Q=H^{-1/2}\dot PH^{-1/2}=\tfrac12\big(H^{-1/2}PH^{-1/2}\big)
\big(H^{-1/2}PH^{-1/2}\big)-2\lambda H^{-1/2}CH^{-1/2}=\tfrac12Q^2-2\lambda
\tilde C,
\]
with $\tilde C=H^{-1/2}CH^{-1/2}=U\,\mathrm{diag}(\mu_i)\,U^\top$. Let
$\hat Q=U^\top QU$; since $U$ is constant,
\[
\dot{\hat Q}=\tfrac12\hat Q^2-2\lambda\,\mathrm{diag}(\mu_i).
\]
To construct the hard-liquidation solution, choose the penalty sequence
$\Phi_a=aH$ with $a\to\infty$. Then $Q(T)=aI$ is diagonal in the $U$
basis. Diagonal matrices are an invariant subspace of the finite-$a$ ODE,
so uniqueness gives diagonal $\hat Q$; taking the limit yields scalar
modes $\dot q_i=\tfrac12q_i^2-2\lambda\mu_i$ with $q_i(T)\to+\infty$.
This argument does not treat infinity as a finite initial condition.

For a scalar Riccati $\dot q=\tfrac12q^2-2\lambda\mu$, substitute
$q=-2\dot\phi/\phi$ (the standard Riccati-to-linear-ODE substitution):
\[
\dot q=-2\Big(\frac{\ddot\phi}{\phi}-\frac{\dot\phi^2}{\phi^2}\Big),\qquad
\tfrac12q^2=2\frac{\dot\phi^2}{\phi^2},
\]
so the Riccati equation becomes $-2\ddot\phi/\phi=-2\lambda\mu$, i.e.\
$\ddot\phi=\lambda\mu\,\phi=\kappa_i^2\phi$ with $\kappa_i=\sqrt{\lambda\mu_i}$.
The blow-up condition $q_i(T)\to\infty$ means $\phi(T)=0$; take
$\phi(t)=\sinh(\kappa_i(T-t))$ (solves $\ddot\phi=\kappa_i^2\phi$ and
vanishes at $t=T$), giving $q_i(t)=2\kappa_i\coth(\kappa_i(T-t))$.

In the decoupled coordinates $\mathbf z(t)=U^\top H^{1/2}\mathbf x(t)$, the
feedback law is
\[
\mathbf v^*=\tfrac12H^{-1}P\mathbf x
=H^{-1/2}U\,\mathrm{diag}(\tfrac12q_i)\,U^\top H^{1/2}\mathbf x.
\]
Together with $\dot{\mathbf x}=-\mathbf v^*$ this gives $\dot z_i=-\kappa_i\coth(\kappa_i(T-t))\,z_i$. Integrating,
$z_i(t)/\sinh(\kappa_i(T-t))=\text{const}$, so $z_i(t)=K_i\sinh(\kappa_i
(T-t))$; matching $z_i(0)=X_i$ (the $i$-th component of $U^\top
H^{1/2}\mathbf X$) fixes $K_i=X_i/\sinh(\kappa_iT)$. Converting back,
$\mathbf x(t)=H^{-1/2}U\mathbf z(t)$, gives the claimed formula.
\end{proof}

For $M=1$ (scalar $H=\eta$, $C=\sigma^2$) this reduces to
$x(t)=X\sinh(\kappa(T-t))/\sinh(\kappa T)$ with
$\kappa=\sqrt{\lambda\sigma^2/\eta}$ --- the continuum limit of Almgren--Chriss
\cite{almgrenchriss2000} Eq.~17. Their finite-step rate instead satisfies
$2(\cosh(\kappa_\tau\tau)-1)/\tau^2=\lambda\sigma^2/\tilde\eta$.
For $\kappa_i=0$, interpret the trajectory ratio as $(T-t)/T$ and
$q_i(t)$ as $2/(T-t)$. For finite $\Phi$, scalar diagonalization requires
$H^{-1/2}\Phi H^{-1/2}$ to commute with $H^{-1/2}CH^{-1/2}$.
Otherwise use the corrected matrix formula in WorkPlan.tex, derived from
Bergault--Drissi--Gu\'eant \cite{bergault2022}, App.~A, Cor.~1.

\section*{Cross-check against the literature}

Every step above is this project's own algebra; the places it is checked
against a published source are: (i) Step~2's continuum limit is the
multi-asset formulation of Bergault--Drissi--Gu\'eant~\cite{bergault2022},
Sec.~2/App.~A, and matches Almgren--Chriss~\cite{almgrenchriss2000},
Sec.~2.2/App.~A.2, in the single-asset case; (ii) Step~B/C's PDE and Riccati
ODE match Bergault--Drissi--Gu\'eant~\cite{bergault2022} Eqs.~16--17
and~20--21 under the stated parameter identification (checked
in the 24 October 2026 audit); (iii) the one-asset trajectory in the
Corollary matches Almgren--Chriss~\cite{almgrenchriss2000} Eq.~17 in the continuum limit.
The symbolic and numerical audit is reproducible with
\texttt{python3 verify\_research.py}; see WorkPlan.tex for the tested
instances, tolerances and remaining limitations.

\begin{thebibliography}{9}

\bibitem{almgrenchriss2000} R.~Almgren and N.~Chriss, ``Optimal Execution of
Portfolio Transactions,'' \emph{J. Risk} \textbf{3}(2), 5--39 (2000).

\bibitem{bergault2022} P.~Bergault, F.~Drissi, and O.~Gu\'eant, ``Multi-asset
Optimal Execution and Statistical Arbitrage Strategies under
Ornstein--Uhlenbeck Dynamics,'' \emph{SIAM J. Financial Math.} \textbf{13}(1),
353--390 (2022); arXiv:2103.13773.

\end{thebibliography}

\end{document}
